\chapter{RELATO DE CASO -- RODOCOCOSE PULMONAR E SALMONELOSE MULTIRRESISTENTE EM POTRO MANGALARGA MARCHADOR}
\label{cap:relato-caso}

\section{Referencial teórico}
\label{sec:referencial}

A rodococose consiste em enfermidade infecciosa de distribuição mundial
causada pela bactéria \textit{Rhodococcus equi}, um microrganismo
gram-positivo, patógeno intracelular facultativo e ubiquitário no solo e
no trato gastrointestinal de herbívoros \cite{Giguere2011, Bordin2022}.
Reconhecida como uma das principais causas de pneumonia grave em
potros com idade entre um e cinco meses de vida, a doença representa
importante desafio para a medicina equina, tanto pela morbidade e
mortalidade associadas quanto pelos elevados custos relacionados ao
tratamento e à vigilância epidemiológica \cite{Cohen2014, Muscatello2012}.

A patogênese da infecção relaciona-se diretamente à capacidade do
\textit{R. equi} de sobreviver e replicar-se no interior de macrófagos
alveolares, mecanismo mediado pelo bloqueio da fusão
fagossomo-lisossomo. Cepas virulentas possuem plasmídeo de virulência
responsável pela codificação da proteína de superfície VapA
(\textit{virulence-associated protein A}), essencial para a sobrevivência
intracelular bacteriana e, consequentemente, para o estabelecimento da
infecção pulmonar \cite{Bordin2022, Lascola2024}. Os potros
geralmente são expostos ao agente nos primeiros dias de vida,
principalmente pela inalação de partículas contaminadas provenientes do
solo, sendo a ingestão também reconhecida como via secundária de
exposição; contudo, os sinais clínicos tendem a manifestar-se entre
quatro e doze semanas de idade \cite{Giguere2011}.

A apresentação clínica clássica da rodococose em potros corresponde à
broncopneumonia piogranulomatosa, frequentemente acompanhada por
abscessos pulmonares multifocais e linfadenite supurativa dos linfonodos
bronquiais e mediastinais \cite{Cohen2014}. Os sinais clínicos iniciais
costumam apresentar caráter inespecífico e insidioso, incluindo febre,
letargia e taquipneia discreta, podendo passar despercebidos nos
estágios iniciais da enfermidade \cite{Cohen2014, Giguere2011}. Com a
progressão do quadro clínico, tornam-se mais evidentes a tosse, a
secreção nasal, o aumento do esforço respiratório e a perda de peso com
retardo no desenvolvimento \cite{Lascola2024, Sanz2023}.

Manifestações extrapulmonares podem ser observadas em até 74\% dos
potros acometidos pela infecção por \textit{R. equi} \cite{Giguere2011}.
Dentre elas, destacam-se o envolvimento intestinal, caracterizado por
enterocolite ulcerativa e linfadenite mesentérica e/ou colônica,
identificado em aproximadamente 50\% dos potros afetados à necropsia,
além da polissinovite imunomediada, observada em cerca de um terço dos
casos. Osteomielite, artrite séptica e abscessos abdominais também podem
ocorrer, embora com menor frequência \cite{Lascola2024}. A sobrevida de
potros com manifestações extrapulmonares mostra-se significativamente
inferior à observada em casos restritos ao comprometimento pulmonar
\cite{Giguere2011}.

O diagnóstico definitivo da rodococose baseia-se no isolamento do
\textit{R. equi} a partir de lavado traqueobrônquico, preferencialmente
associado à identificação molecular do gene \textit{vapA} por reação em
cadeia da polimerase (PCR) \cite{Cohen2014, Muscatello2012}. A
ultrassonografia torácica representa a principal ferramenta para
detecção precoce de consolidações pulmonares e abscessos, enquanto a
radiografia torácica complementa a avaliação por meio da identificação
de padrões alveolares e intersticiais, consolidação regional e formações
abscedativas \cite{Cohen2014, Muscatello2012}. Resultados negativos em
culturas obtidas por swab nasal ou lavado traqueal não excluem a
infecção, uma vez que o agente pode não ser isolado em todas as amostras
obtidas por essas vias \cite{Sanz2023}.

O tratamento padrão da rodococose em potros consiste na associação entre
um antimicrobiano macrolídeo, como azitromicina ou claritromicina, e
rifampicina, administrados por período prolongado \cite{Bordin2022, Cohen2014}.
Essa associação terapêutica fundamenta-se na elevada
lipossolubilidade de ambas as classes farmacológicas, favorecendo a
penetração em piogranulomas e o alcance intracelular das bactérias.
Entretanto, nos últimos anos, tem sido observada emergência crescente de
cepas multirresistentes de \textit{R. equi}, especialmente em propriedades
nas quais ocorre uso profilático indiscriminado com macrolídeos e
rifampicina \cite{Bordin2022, Giguere2017}. Tal cenário limita as
opções terapêuticas disponíveis e representa preocupação relevante tanto
para a saúde equina quanto para a saúde pública.

A salmonelose figura entre as enfermidades infecciosas de maior
relevância na medicina equina, sendo causada por diferentes sorovares de
\textit{Salmonella enterica}. A doença pode acometer equinos de diferentes
faixas etárias; contudo, potros apresentam maior susceptibilidade em
razão da imaturidade imunológica, o que favorece o desenvolvimento de
bacteremia, septicemia e infecções articulares secundárias
\cite{TraubDargatz2016}.

As manifestações clínicas da salmonelose em equinos incluem diarreia
aquosa profusa, febre, cólica, desidratação e sinais de septicemia. Nas
formas peragudas, o animal pode evoluir rapidamente para óbito, mesmo
diante de terapia intensiva. A enteropatia perdedora de proteínas pode
resultar em hipoproteinemia grave, enquanto a translocação bacteriana
intestinal pode favorecer bacteremia e coagulação intravascular
disseminada \cite{MerckVet2025}.

O tratamento antimicrobiano da salmonelose em equinos adultos permanece
controverso, sendo geralmente reservado para animais com neutropenia
grave, febre persistente ou utilização de cateter intravenoso. Em
potros, entretanto, a antibioticoterapia mostra-se indicada devido à
elevada predisposição ao desenvolvimento de septicemia. Os
antimicrobianos mais empregados incluem ceftiofur, enrofloxacina e
aminoglicosídeos \cite{Cummings2016}.

A multirresistência antimicrobiana em isolados de \textit{Salmonella} de
origem equina constitui preocupação crescente e amplamente documentada
na literatura. Estudos demonstram que mais de 50\% dos isolados de
\textit{Salmonella} provenientes de equinos hospitalizados apresentam
perfil de multirresistência \cite{Cummings2016}. Essa condição reduz
significativamente as possibilidades terapêuticas e aumenta o risco de
falha no tratamento, comprometendo o prognóstico dos animais acometidos.

No ambiente hospitalar veterinário, \textit{Salmonella enterica}
representa agente de importante relevância nosocomial. Surtos de
salmonelose em hospitais veterinários equinos encontram-se amplamente
descritos na literatura, sendo caracterizados por elevadas taxas de
morbidade e mortalidade, transmissão nosocomial entre pacientes,
potencial zoonótico para profissionais e, em situações extremas,
necessidade de interrupção temporária das atividades hospitalares
\cite{SozaOssandon2020, Schott2001}. A persistência
ambiental de cepas multirresistentes de \textit{Salmonella} em hospitais
veterinários constitui desafio adicional para o controle da infecção,
reforçando a necessidade de protocolos rigorosos de biossegurança,
vigilância epidemiológica e uso racional de antimicrobianos
\cite{SozaOssandon2020}.

A coinfecção por \textit{R. equi} e \textit{Salmonella} spp. em potros,
embora raramente descrita como entidade clínica isolada, configura
cenário de elevada gravidade clínica. A associação entre broncopneumonia
abscedativa, enterite e septicemia causada por \textit{Salmonella}
multirresistente resulta em quadro de elevada complexidade terapêutica,
comprometendo significativamente o prognóstico do paciente.

\section{Relato de caso}
\label{sec:relato}

Um potro da raça Mangalarga Marchador, com dois meses de vida e pesando
75 kg, foi encaminhado ao setor de Clínica Médica de Equinos do Hospital
Veterinário da Universidade Federal de Minas Gerais (HV-UFMG) após o
proprietário observar aumento de volume bilateral em região de jarretes.
A suspeita inicial de artrite séptica, levantada ainda na propriedade
pela médica veterinária responsável pelo primeiro atendimento, motivou a
instituição de tratamento prévio com ceftiofur, amicacina e meloxicam,
sendo o animal posteriormente encaminhado ao hospital diante da
gravidade do quadro.

Na manhã do dia anterior em que foi encaminhado ao hospital, o quadro se
agravou com manifestação de apatia, recusa alimentar e hipertermia.

No dia 18 de abril de 2026, ao atendimento inicial no setor, o paciente
encontrava-se apático e alheio ao ambiente. Ao exame físico,
observaram-se frequência cardíaca de 130 bpm, pulso periférico fraco,
frequência respiratória de 57 movimentos por minuto com padrão
costoabdominal acentuado e estertores evidentes, mais perceptíveis nos
focos laringotraqueal e traqueobrônquico. As mucosas apresentavam-se
hipocoradas, com tempo de preenchimento capilar de dois segundos e
temperatura retal de 39,3~°C. Os parâmetros gastrointestinais não
apresentaram alterações relevantes e a glicemia aferida foi de 111
mg/dL.

Diante do comprometimento respiratório, realizou-se ultrassonografia
pulmonar que evidenciou múltiplas áreas hipoecoicas, bem delimitadas,
compatíveis com consolidação pulmonar e formação de abscessos,
associadas à irregularidade pleural e à presença de linhas B, as
chamadas ``caudas de cometa'', indicativas de processo inflamatório
pulmonar multifocal. A radiografia torácica complementou os achados,
demonstrando aumento multifocal da radiopacidade em padrão intersticial
difuso, compatível com broncopneumonia extensa.

\begin{figure}[htb]
  \centering
  \caption{Imagens da ultrassonografia pulmonar.}
  \label{fig:usg-pulmonar}
  % [INSERIR IMAGEM]
  \fonte{Setor de Clínica Médica de Equinos/UFMG (2026).}
\end{figure}

\begin{figure}[htb]
  \centering
  \caption{Radiografia torácica.}
  \label{fig:rx-toracica}
  % [INSERIR IMAGEM]
  \fonte{Setor de Clínica Médica de Equinos/UFMG (2026).}
\end{figure}

Paralelamente, procedeu-se à avaliação ortopédica das articulações,
motivada pela queixa inicial de aumento de volume bilateral nos
jarretes. Ao exame físico, constatou-se que o aumento de volume
estendia-se também aos boletos bilateralmente. À ultrassonografia,
observou-se distensão da cápsula articular com acúmulo de conteúdo
anecoico, compatível com efusão sinovial. A radiografia não revelou lise
óssea nem proliferação osteoarticular, achado que, embora tranquilizador
em relação à integridade estrutural, não descartava comprometimento
infeccioso articular. A análise macroscópica do líquido sinovial revelou
coloração amarela clara, dentro do esperado. À análise laboratorial,
demonstrou alterações inflamatórias leves a moderadas, baixa
viscosidade, discreta turbidez e predomínio neutrofílico, porém sem as
alterações clássicas de artrite séptica franca, como contagem celular
extremamente elevada ou aspecto purulento.

\begin{figure}[htb]
  \centering
  \caption{Imagem do líquido sinovial coletado.}
  \label{fig:liquido-sinovial}
  \includegraphics[width=0.5\textwidth]{media/image12.jpg}
  \fonte{Setor de Clínica Médica de Equinos/UFMG (2026).}
\end{figure}

\begin{figure}[htb]
  \centering
  \caption{Imagens da ultrassonografia das articulações do tarso do membro pélvico esquerdo.}
  \label{fig:usg-tarso}
  % [INSERIR IMAGEM]
  \fonte{Setor de Clínica Médica de Equinos/UFMG (2026).}
\end{figure}

\begin{figure}[htb]
  \centering
  \caption{Imagens da radiografia das articulações do carpo do membro torácico esquerdo.}
  \label{fig:rx-carpo}
  % [INSERIR IMAGEM]
  \fonte{Setor de Clínica Médica de Equinos/UFMG (2026).}
\end{figure}

O hemograma revelou anemia, com hematócrito 16\%, desvio à esquerda dos
neutrófilos e presença de neutrófilos tóxicos. No perfil bioquímico
sérico, verificou-se aumento da ureia (63,24 mg/dL) compatível com
catabolismo proteico e desidratação, redução das proteínas totais (4,81
g/dL) e globulinas (2,4 g/dL), sugestiva de perda ou déficit de síntese
proteica, possivelmente associada ao comprometimento sistêmico grave, e
hipertrigliceridemia (481,25 mg/dL), refletindo desequilíbrio energético
e inflamação sistêmica.

\begin{figure}[htb]
  \centering
  \caption{Resultado do eritrograma do paciente.}
  \label{fig:eritrograma}
  % [INSERIR IMAGEM]
  \fonte{Setor de Clínica Médica de Equinos/UFMG (2026).}
\end{figure}

\begin{figure}[htb]
  \centering
  \caption{Resultado do leucograma do paciente.}
  \label{fig:leucograma}
  % [INSERIR IMAGEM]
  \fonte{Setor de Clínica Médica de Equinos/UFMG (2026).}
\end{figure}

\begin{figure}[htb]
  \centering
  \caption{Resultado do perfil bioquímico sérico do paciente.}
  \label{fig:bioquimico}
  % [INSERIR IMAGEM]
  \fonte{Setor de Clínica Médica de Equinos/UFMG (2026).}
\end{figure}

A hemogasometria trouxe resultados compatíveis com acidose metabólica
importante, caracterizada por redução do pH sanguíneo (7,191),
diminuição do bicarbonato sérico (8,8 mmol/L) e déficit de base
(-17,7 mmol/L), com aumento do ânion gap (32,2 mmol/L) e elevação do
lactato sérico (5,4 mmol/L) indicativos de acúmulo de ácidos circulantes
e comprometimento da perfusão tecidual. Hipocalemia (2,6 mEq/L)
completava o painel de distúrbios eletrolíticos.

\begin{figure}[htb]
  \centering
  \caption{Compilado das alterações identificadas na hemogasometria do paciente.}
  \label{fig:hemogasometria}
  % [INSERIR IMAGEM]
  \fonte{Da autora (2026).}
\end{figure}

Entre as suspeitas etiológicas consideradas, o \textit{Rhodococcus equi}
figurou em primeiro plano, dada a combinação de broncopneumonia com
formação de abscessos pulmonares em um potro jovem. Foram realizados
swab nasal e lavado traqueal para pesquisa do agente, com resultados
negativos em ambas as amostras. Durante a investigação, a pesquisa
molecular em swab nasal detectou positividade para \textit{Equine
Herpesvirus 5} (EHV-5); entretanto, considerando os achados clínicos,
ultrassonográficos, radiográficos e os achados anatomopatológicos
obtidos posteriormente, esse resultado foi interpretado como um achado
incidental, sem relação causal com o quadro clínico apresentado pelo
paciente.

O tratamento foi instituído de forma intensiva. A fluidoterapia
intravenosa com Ringer com lactato visou corrigir a desidratação e
restaurar o equilíbrio hemodinâmico. Com base nos resultados
hemogasométricos, procedeu-se à suplementação de bicarbonato para
correção da acidose metabólica e à reposição de potássio. O protocolo
antimicrobiano manteve ceftiofur (10 mg/kg, IV, BID) e amicacina (25
mg/kg, IV, SID), já iniciados antes da internação. Para controle da
inflamação sistêmica e analgesia, administrou-se flunixin meglumine (1,1
mg/kg, IV, BID). O suporte respiratório foi conduzido com nebulizações
realizadas duas vezes ao dia contendo acetilcisteína (1,5 mL),
bromexina (1,5 mL) e solução fisiológica (2 mL), visando a fluidificação
das secreções e facilitação da drenagem das vias aéreas.

Para minimizar os riscos de lesão gástrica associada ao uso de
anti-inflamatório não esteroidal, utilizou-se gastrozol (pasta para 75
kg, VO, SID). A preocupação com a microbiota intestinal levou à inclusão
de probiótico Pró-Sacc no protocolo (VO, SID), antecipando o risco de
disbiose decorrente do uso prolongado de antimicrobianos. Diante da
dificuldade de amamentação espontânea, a égua mãe foi ordenhada
manualmente e o leite fornecido em balde ao potro, que demonstrou boa
aceitação.

Ao longo da internação, observou-se melhora progressiva do padrão
respiratório e da ausculta pulmonar. O potro, no entanto, apresentou
episódios de diarreia fétida de coloração amarelo-esverdeada, sendo
investigada salmonelose. O isolamento bacteriano confirmou a infecção
por \textit{Salmonella spp.} com perfil de resistência antimicrobiana a
múltiplos agentes, o que limitou significativamente as opções
terapêuticas disponíveis e agravou substancialmente o manejo clínico. As
alterações hematológicas persistiram, com leucocitose mantida e
variações no hematócrito, mesmo com transfusões, que evidenciaram
ausência de controle efetivo do processo infeccioso sistêmico.

Diante do comprometimento pulmonar multifocal extenso associado a
salmonelose multirresistente sem opções terapêuticas viáveis e piora
sistêmica progressiva, o prognóstico foi considerado desfavorável. A
equipe, em conjunto com o responsável pelo animal, optou pela realização
de eutanásia.

À necropsia, foram observadas múltiplas áreas branco-amareladas,
multifocais, no parênquima pulmonar, que ao corte drenavam conteúdo
purulento, achado macroscópico clássico da rodococose pulmonar. Também
foram evidenciados aumento da cadeia linfonodal do cólon e incremento do
volume do líquido sinovial nas articulações avaliadas, indicando
disseminação sistêmica do processo infeccioso. O exame anatomopatológico
posterior confirmou a infecção por \textit{Rhodococcus equi}, corroborando
a suspeita clínica levantada desde a admissão, a despeito dos resultados
negativos obtidos nos exames \textit{ante mortem}.

\begin{figure}[htb]
  \centering
  \caption{Achados macroscópicos observados na necropsia do paciente.}
  \label{fig:necropsia}
  % [INSERIR IMAGEM]
  \fonte{Setor de Clínica Médica de Equinos/UFMG (2026).}
\end{figure}

\section{Discussão}
\label{sec:discussao}

O caso relatado ilustra a complexidade diagnóstica e terapêutica da
rodococose pulmonar em potros, especialmente quando associada à
coinfecção por \textit{Salmonella} spp. multirresistente. O paciente, um
potro da raça Mangalarga Marchador com dois meses de idade, apresentou
quadro clínico compatível com a literatura, a qual descreve maior
susceptibilidade à infecção por \textit{Rhodococcus equi} em animais entre
um e cinco meses de vida \cite{Cohen2014, Lascola2024}.

A apresentação clínica observada, caracterizada por taquipneia,
taquicardia, padrão respiratório costoabdominal acentuado, estertores
pulmonares e febre, mostrou-se compatível com os achados descritos em
casos avançados de rodococose, nos quais a evolução insidiosa da
enfermidade frequentemente retarda o diagnóstico até estágios mais
graves \cite{Giguere2011, Sanz2023}. A identificação ultrassonográfica de
múltiplas áreas hipoecoicas compatíveis com consolidação pulmonar e
abscessos, posteriormente confirmadas pela radiografia torácica, reforça
a importância dos exames de imagem como ferramentas fundamentais para o
diagnóstico da rodococose \cite{Cohen2014, Muscatello2012}.

A distensão articular bilateral observada nos jarretes, motivo inicial
do encaminhamento do animal, mostrou-se compatível com a polissinovite
imunomediada descrita em aproximadamente um terço dos potros acometidos
por infecção por \textit{R. equi}. Essa manifestação extrapulmonar decorre
da deposição de imunocomplexos em estruturas sinoviais e, classicamente,
ocorre na ausência de claudicação, com líquido sinovial apresentando
pleocitose mononuclear sem crescimento bacteriano \cite{Giguere2011, Lascola2024}.
Os achados obtidos na análise do líquido sinovial do
paciente, caracterizados por alterações inflamatórias leves a moderadas
sem evidências compatíveis com artrite séptica, corroboram essa
interpretação.

Os resultados negativos para \textit{R. equi} nos exames de swab nasal e
lavado traqueal merecem destaque, uma vez que o diagnóstico definitivo
da rodococose fundamenta-se preferencialmente no isolamento do agente a
partir de aspirado traqueobrônquico \cite{Cohen2014}. A limitada
sensibilidade dessas técnicas de coleta, associada à possível
interferência do tratamento antimicrobiano previamente instituído na
propriedade, com ceftiofur e amicacina, pode justificar os resultados
falso-negativos observados. A confirmação diagnóstica foi obtida apenas
no exame post-mortem, mediante identificação de lesões
anatomopatológicas características da rodococose, incluindo abscessos
pulmonares multifocais com conteúdo purulento e linfadenite colônica,
posteriormente confirmadas pelo exame histopatológico.

As alterações laboratoriais identificadas, incluindo anemia grave,
desvio à esquerda associado à presença de neutrófilos tóxicos,
hipoproteinemia, hipertrigliceridemia e acidose metabólica, refletem a
gravidade do processo infeccioso sistêmico e são consistentes com o
perfil descrito na rodococose avançada. O desvio à esquerda com
neutrofilia tóxica, em particular, é concordante com o padrão
leucocitário descrito na literatura, no qual são frequentemente
relatadas leucocitose neutrofílica e hiperfibrinogenemia como marcadores
do processo inflamatório ativo \cite{Cohen2014, Lascola2024}.

O desenvolvimento de diarreia fétida de coloração amarelo-esverdeada
durante o período de internação, com posterior confirmação de infecção
por \textit{Salmonella} spp. multirresistente, caracterizou um quadro de
salmonelose nosocomial, uma vez que o animal foi admitido sem sinais
gastrointestinais prévios. Surtos de salmonelose em hospitais
veterinários equinos estão amplamente documentados na literatura, sendo
\textit{Salmonella enterica} reconhecida como um dos principais agentes de
infecção hospitalar nesse contexto, com elevado potencial de transmissão
entre pacientes e risco zoonótico para os profissionais envolvidos
\cite{SozaOssandon2020, Schott2001}.

O perfil de multirresistência antimicrobiana identificado no isolado
reduziu criticamente as opções terapêuticas disponíveis, corroborando os
achados descritos por \citeonline{Cummings2016}, que identificaram
multirresistência em mais da metade dos isolados de \textit{Salmonella}
provenientes de equinos hospitalizados. Nesse sentido, o uso prévio de
ceftiofur e amicacina, instituído antes da internação, pode ter exercido
pressão seletiva sobre cepas resistentes, contribuindo para o perfil
encontrado \cite{Bordin2022}. A coinfecção por \textit{R. equi} e
\textit{Salmonella} multirresistente resultou em quadro de complexidade
terapêutica extrema, no qual a falha no controle de ambos os processos
infecciosos simultaneamente inviabilizou resposta clínica satisfatória,
tornando a eutanásia a conduta mais adequada sob a perspectiva do
bem-estar animal.

O caso evidencia a gravidade que a rodococose pode alcançar quando o
diagnóstico é tardio, reforçando a importância dos exames de imagem e da
investigação clínica criteriosa como ferramentas essenciais para o
reconhecimento precoce da doença. Paralelamente, a multirresistência
identificada no isolado de \textit{Salmonella} spp. reforça a necessidade
de protocolos rigorosos de biossegurança, vigilância epidemiológica e
uso racional de antimicrobianos no ambiente hospitalar veterinário, bem
como o desenvolvimento de novas estratégias profiláticas frente à
crescente emergência de cepas multirresistentes de \textit{R. equi} em
nível mundial \cite{Bordin2022, Giguere2017}.
